% SOLUCIÓN Ejercicio 1
\documentclass[12pt,a4paper]{article}
\usepackage[utf8]{inputenc}
\usepackage[T1]{fontenc}
\usepackage[spanish]{babel}
\usepackage{amsmath,amssymb}
\usepackage{booktabs}
\usepackage{siunitx}
\sisetup{locale=FR, output-decimal-marker={,}}
\usepackage{circuitikz}
\usepackage[hidelinks]{hyperref}

\title{Ecuaciones de electrónica}
\author{Estudiante UNTELS}
\date{\today}

\begin{document}
\maketitle

\section{Ecuaciones fundamentales}
\begin{enumerate}
  \item Ley de Ohm:
    \begin{equation} V = I\cdot R \label{eq:ohm} \end{equation}
  \item Potencia eléctrica:
    \begin{equation} P = V I = I^{2}R = \frac{V^{2}}{R} \label{eq:potencia} \end{equation}
  \item Reactancia capacitiva:
    \begin{equation} X_C = \frac{1}{2\pi f C} \label{eq:xc} \end{equation}
  \item Reactancia inductiva:
    \begin{equation} X_L = 2\pi f L \label{eq:xl} \end{equation}
  \item Impedancia RLC serie:
    \begin{equation} Z = \sqrt{R^{2} + (X_L - X_C)^{2}} \label{eq:z} \end{equation}
  \item Frecuencia de resonancia:
    \begin{equation} f_0 = \frac{1}{2\pi\sqrt{LC}} \label{eq:f0} \end{equation}
\end{enumerate}

\section{Cálculo}
Con $R=\SI{100}{\ohm}$, $L=\SI{10}{\milli\henry}$ y $C=\SI{1}{\micro\farad}$,
la ecuación~\eqref{eq:f0} da
\[
  f_0 = \frac{1}{2\pi\sqrt{10\times10^{-3}\cdot 1\times10^{-6}}}
      \approx \SI{1,59}{\kilo\hertz}.
\]
En resonancia $X_L=X_C$ y, por~\eqref{eq:z}, $Z=R=\SI{100}{\ohm}$.

\end{document}
